\chapter{Diskussion}\label{chap:discussion}

  \section{Interpretation der Ergebnisse}\label{sec:interpretation-of-results}
    Das primäre Ziel dieser Arbeit bestand in der prototypischen Entwicklung eines 6-Achsen-Slicers für die nicht-planare, robotergestützte additive Fertigung.
    Die vorliegende Arbeit adressiert dabei die signifikante Diskrepanz zwischen der bereits hoch entwickelten 6-Achs-Hardware und der bisher unzureichenden Verfügbarkeit entsprechender Standard-Slicing-Software.
    Im Rahmen der Implementierung ist es erfolgreich gelungen, eine durchgängige und konsistente Datenkette zu etablieren, die von der Verarbeitung roher \ac{STL}-Daten bis hin zur präzisen Generierung komplexer Roboterbewegungen reicht.
    Damit konnte die grundlegende Machbarkeit eines solchen Systems sowie dessen technologisches Potenzial zweifelsfrei nachgewiesen werden.

    Die explizite Trennung zwischen \ac{medusa} für das Slicing und \ac{kronos} für die Robotersteuerung verdeutlicht eine modulare Systemarchitektur, die eine hohe Flexibilität für zukünftige Anpassungen bietet.
    Dass der Datentransfer via \ac{JSON} aktuell noch manuell erfolgt, markiert zwar eine Unterbrechung im automatisierten Workflow, unterstreicht jedoch gleichzeitig die Rolle von \ac{medusa} als eigenständiges, plattformunabhängiges Werkzeug zur Pfadgenerierung.
    Eine zukünftige Automatisierung dieses Schrittes könnte die Effizienz weiter steigern, ohne die grundsätzliche Funktionalität zu beeinträchtigen.
    Dies könnte beispielsweise über eine \ac{REST}-\ac{API} oder Pipes realisiert werden.
    Die erfolgreiche Validierung der Trajektorien innerhalb der \ac{kronos}-Simulation ist als kritischer Sicherheitsfaktor zu interpretieren, da sie das Risiko von Hardware-Kollisionen im realen Betrieb des UR3e signifikant minimiert.

    Obgleich an spezifischen Stellen des Systems noch Optimierungspotenzial sowie ein Bedarf an weiterführender Entwicklung identifiziert wurde, erfüllt dieser Prototyp seine Funktion als Proof-of-Concept vollumfänglich.
    Die Ergebnisse zeigen auf, dass das gewählte Verfahren operiert und eine belastbare Basis für zukünftige Forschungsvorhaben darstellt.
    Somit legt diese Arbeit den entscheidenden Grundstein für eine solide Weiterentwicklung hin zu einem produktionsreifen System in der adaptiven Fertigung.
    \sig{Ben Stahl}

    \section{Entwurfsentscheidungen} \label{sec:discussion-design-decisions}
      Die softwaretechnische Implementierung des Systems \ac{medusa} basiert auf einer strikt geschichteten Modularchitektur, welche konsequent auf dem Einsatz statischer Bibliotheken und einheitlicher Namespace-Aliase aufbaut.
      Diese architektonische Struktur hat sich im Projektverlauf als äußerst tragfähig und skalierbar erwiesen.
      Ein wesentlicher Erfolgsfaktor ist hierbei die konsequente Trennung der fachlichen Domänen zwischen der geometrischen Pfadplanung, die durch das Modul \ac{medusa} repräsentiert wird, und der roboterkinematischen Ausführung, welche durch das Teilsystem \ac{kronos} realisiert wird.

      Während das Teilsystem \ac{medusa} primär High-Level-Aufgaben wie den fehlerfreien Mesh-Import von Geometriedaten sowie die komplexen Slicing-Operationen übernimmt, fokussiert sich \ac{kronos} dediziert auf die Berechnung der Inversen Kinematik, die dynamische Kollisionsprüfung und die hardwarenahe Ansteuerung des Manipulators.
      Diese klare funktionale Aufteilung erlaubt nicht nur den isolierten und zielgerichteten Einsatz fachspezifischer Programmbibliotheken, sondern trägt auch den grundlegend unterschiedlichen Laufzeitanforderungen der beiden Systeme Rechnung.
      \ac{medusa} agiert hierbei als reine Offline-Anwendung ohne strikte Echtzeitbedingungen, wohingegen \ac{kronos} inhärent an die Kommunikationszyklen und die hochkritischen Sicherheitsüberwachungen der realen Robotersteuerung gebunden ist.

      Ein wesentliches technisches Detail zur Sicherung der Softwarequalität innerhalb von \ac{medusa} betrifft die Definition einer strikten „Include-only“-Abhängigkeit zwischen den Modulen \texttt{graphics} und \texttt{slicing}.
      Durch diese gezielte Entkopplung werden zirkuläre Abhängigkeiten im Quellcode effektiv vermieden, was die Wartbarkeit des Systems erheblich steigert.
      Gleichwohl erfordert diese restriktive Architektur eine hohe disziplinäre Sorgfalt bei zukünftigen Erweiterungen des \texttt{ToolpathRenderer}, um die softwareseitige Integrität der Visualisierungsschicht langfristig zu gewährleisten.
      \sig{Ben Stahl}

  \section{Validität und Übertragbarkeit} \label{sec:validity-and-transferability}
    \subsection{Interne und externe Validität} \label{subsec:internal-and-external-validity}
    Die interne Validität des Systems wird maßgeblich durch die Simulationsumgebung innerhalb des Teilsystems \ac{kronos} gewährleistet, welche eine Überprüfung der von \ac{medusa} generierten Toolpaths ermöglicht.
    Durch diese virtuelle Validierung können die berechneten Trajektorien auf ihre geometrische Korrektheit und algorithmische Konsistenz hin untersucht werden.
    Kritisch zu berücksichtigen bleibt jedoch die in Abschnitt \ref{subsec:limitations-and-failure-scenarios} detailliert beschriebene Problematik hinsichtlich des Inverse-Kinematik-Solvers, welche die Zuverlässigkeit der Ergebnisse in spezifischen Randbereichen beeinflussen kann.

    Hinsichtlich der externen Validität muss festgehalten werden, dass eine vollumfängliche Überprüfung unter realen Produktionsbedingungen zum aktuellen Zeitpunkt noch aussteht.
    Dies liegt primär darin begründet, dass eine Validierung mittels realer Materialextrusion oder weiterführender physikalischer Belastungstests nicht im ursprünglichen Scope der Projektziele definiert wurde.
    Die externe Validität beschränkt sich somit planmäßig auf die digitale Repräsentation und die Verifikation der theoretischen Bewegungsabläufe, während die physische Umsetzung ein Feld für zukünftige Forschungsarbeiten darstellt.
    
    Ein weiterer wesentlicher Aspekt im Hinblick auf die Validität der vorliegenden Ergebnisse ist der Umstand, dass zum aktuellen Zeitpunkt noch kein expliziter Materialfluss innerhalb der Simulation abgebildet wird.
    Infolgedessen lässt sich das Verhalten des extrudierten Materials während des eigentlichen Schichtaufbaus nicht mit absoluter Gewissheit vollumfänglich prognostizieren.
    Da die physikalischen Interaktionen, wie Gravitation, beim Austritt aus der Düse eine entscheidende Rolle für die strukturelle Integrität des Bauteils spielen, bleibt die Vorhersage der Materialdynamik ungewiss.
    Mit dieser Thematik wird sich in Abschnitt \ref{subsec:sideways-material-extrusion} genauer befasst.
    \sig{Ben Stahl}
    
    \subsection{Reproduzierbarkeit der Ergebnisse}\label{subsec:reproducibility-of-results}
      Die Reproduzierbarkeit der generierten Roboterbewegungen des Universal Robots UR3e wird primär durch den Einsatz von RViz innerhalb der \ac{ROS} 2-Umgebung sichergestellt.
      Diese Simulationsplattform ermöglicht es, die berechneten Trajektorien unter identischen Randbedingungen wiederholt zu visualisieren und deren Ausführung zu verifizieren.
      Ein wesentlicher technischer Aspekt, der die absolute Reproduzierbarkeit jedoch einschränkt, ist die Natur des verwendeten Inverse-Kinematik-Solvers.

      Da die Berechnung der Gelenkwinkel für eine vorgegebene Pose im Raum auf numerischen Iterationsverfahren basiert, ist die Lösungsfindung nicht in jedem Fall streng deterministisch.
      Dies bedeutet, dass für identische Zielkoordinaten potenziell unterschiedliche Konfigurationen der Roboterachsen berechnet werden können, sofern mehrere mathematisch valide Lösungen existieren.
      Infolgedessen ist eine Reproduzierbarkeit der Bewegungsabläufe zwar grundsätzlich gegeben, ihre absolute Zuverlässigkeit muss jedoch im Kontext der algorithmischen Varianz des Solvers kritisch betrachtet werden.
      \sig{Ben Stahl}

  \section{Limitationen der Arbeit}\label{sec:thesis-limitations}
    \subsection{Methodische Limitationen}\label{subsec:methodological-limitations}
      Zu den wesentlichen methodischen Limitationen dieses Projektes zählt primär der Verzicht auf die Implementierung des sphärenbasierten Algorithmus.
      Wie bereits in der detaillierten Evaluation in Abschnitt \ref{subsec:evaluation-of-sphere-slicer} dargelegt wurde, erwies sich dieser Ansatz aufgrund struktureller Inkonsistenzen in Überhangbereichen als suboptimal für die angestrebten Fertigungsziele.
      Jedoch wurde, wie angemerkt, die weiterentwickelte verbesserte Version implementiert.

      Ein weiterer limitierender Faktor im aktuellen Systemaufbau ist das Fehlen einer dynamischen Kollisionsprüfung zwischen dem Manipulator des Universal Robots UR3e und dem bereits extrudierten Material.
      Es muss jedoch hervorgehoben werden, dass diese spezifische Funktionalität gemäß der Anforderungsdefinition in Abschnitt \ref{subsec:kronos-requirements} explizit nicht Bestandteil des aktuellen Projektumfangs war.

      Für eine zukünftige Weiterführung des Projektes, die insbesondere auf die reale physische Extrusion abzielt, stellt dieser Aspekt jedoch ein kritisches Entwicklungsfeld dar.
      Um eine prozesssichere Fertigung komplexer Geometrien zu gewährleisten, ist die Integration einer Echtzeit-Kollisionsüberwachung für das gedruckte Objekt zwingend erforderlich.
      Dies erfordert zusätzliche Entwicklungsarbeit im Bereich der Pfadplanung und der dynamischen Umfeldwahrnehmung innerhalb der Simulationsumgebung.
      \sig{Ben Stahl}

    \subsection{Technische Limitationen}\label{subsec:technical-limitations}
      Die technischen Limitationen dieses Projektes manifestieren sich primär in den in Abschnitt \ref{subsec:limitations-and-failure-scenarios} detailliert erörterten Problematiken hinsichtlich der Berechnung nicht-linearer Pfadsegmente.
      Ein zentraler Schwachpunkt liegt dabei in der mangelnden Linearität der Trajektorien zwischen definierten Wegpunkten, was zu unerwünschten Abweichungen in der angestrebten Bahnbewegung führen kann.
      Des Weiteren stellt das gelegentliche Ausbleiben einer validen Bewegungslösung durch den Inverse-Kinematik-Solver eine signifikante Hürde für die Prozessstabilität dar.

      Für jede zukünftige Weiterentwicklung oder gar eine Überführung des Systems in den produktiven Einsatz ist die Behebung dieser Defizite von essenzieller Bedeutung.
      Nur durch eine Optimierung der kinematischen Algorithmen und eine robustere Pfadplanungslogik kann eine zuverlässige und präzise Ansteuerung des Universal Robots UR3e dauerhaft gewährleistet werden.
      \sig{Ben Stahl}

  \section{Erweiterbarkeit} \label{sec:discussion-extensibility}
    Die Softwarearchitektur von \ac{medusa} wurde gezielt auf eine hohe strukturelle Flexibilität und langfristige Erweiterbarkeit ausgelegt.
    Durch die konsequente Implementierung des \texttt{ISlicer}-Interfaces wird eine abstrakte Schnittstelle geschaffen, welche die Integration neuartiger Slicing-Algorithmen maßgeblich vereinfacht.
    Entwickler können somit maßgeschneiderte Slicing-Strategien implementieren und nahtlos in das bestehende System einbetten, ohne dass Modifikationen an der übergeordneten Pipeline-Logik oder an nachgelagerten Verarbeitungsstufen vorgenommen werden müssen.
    Diese lose Kopplung minimiert das Risiko von Regressionseffekten bei Systemerweiterungen und sichert die funktionale Integrität des Systems.

    Ein weiterer architektonischer Vorteil manifestiert sich in der Definition der Projektstruktur auf Top-Level-Ebene.
    Diese saubere Konfiguration erlaubt es, das gesamte \ac{medusa}-Framework ohne tiefgreifenden Anpassungsaufwand als eigenständiges \texttt{CMake}-Unterprojekt in übergeordnete Softwaresysteme oder industrielle Toolchains zu integrieren.
    Auch die technologische Portabilität im Hinblick auf die verwendete Hardware wurde im Designprozess antizipiert.
    Für die Anbindung des Slicers an alternative, herstellerfremde Robotersysteme wäre lediglich die Implementierung eines zusätzlichen, isolierten Moduls erforderlich.

    Dieses dedizierte Modul müsste ausschließlich die mathematische Transformation der generierten Segment-Orientierungen in die spezifischen 6-\ac{DOF}-Posen des jeweiligen Manipulators kapseln.
    Gleichwohl muss kritisch angemerkt werden, dass das Teilsystem \ac{kronos} in seiner aktuellen Implementierungsstufe spezifisch auf die Kinematik und die Parameter des Universal Robots UR3e ausgelegt ist.
    Für die erfolgreiche Manipulation und Ansteuerung anderer Robotertypen ist daher eine entsprechende softwareseitige Erweiterung der kinematischen Modelle innerhalb von \ac{kronos} zwingend erforderlich.
    Die konsequente Nutzung des Frameworks \ac{ROS} 2 als zugrundeliegendes Operating System, welches bereits für die Steuerung des UR3e eingesetzt wird, erleichtert diesen Migrationsprozess jedoch erheblich.
    Durch die inhärente Modularität und die standardisierten Hardware-Interfaces von \ac{ROS} 2 über MoveIt 2 können alternative Manipulatoren mit vergleichsweise geringem Integrationsaufwand eingebunden werden.
    Diese technologische Flexibilität unterstreicht die universelle Einsetzbarkeit und die zukunftssichere Ausrichtung der entwickelten Systemarchitektur für die robotergestützte Fertigung.
    \sig{Ben Stahl}

  \section{Diskussion der Forschungsfragen}
    \subsection{Forschungsfrage 1}
      \subsubsection{Welche Anforderungen muss ein Slicer erfüllen, um nicht-planare Bahnstrategien für einen 6-Achs-Roboter zu erzeugen?}
        Um die zusätzlichen Freiheitsgrade eines 6-Achs-Roboters vollständig zu nutzen, muss ein Slicer weitaus komplexere Aufgaben bewältigen als ein herkömmlicher planarer Slicer.
        Eine zentrale Anforderung ist die simultane Planung von Werkzeugbahn (Toolpath), Werkzeugorientierung und Bewegungsablauf.
        Dies vergrößert den Lösungsraum und die Komplexität der einzuhaltenden Randbedingungen erheblich.
        Wie in \ref{subsec:technical-challenges-path-planning} dargelegt, müssen zudem der Materialfluss aus der Düse, die Bahngeschwindigkeit und die Orientierung des Werkzeugs präzise aufeinander abgestimmt werden, um einen gleichmäßigen Materialauftrag sicherzustellen.

        Hierfür muss der Slicer für jeden Wegpunkt eine vollständige 6-\ac{DOF}-Pose bereitstellen.
        Die Orientierung wird dabei idealerweise in Form von Einheits-Quaternionen an die Robotersteuerung übergeben, um Singularitäten, wie den Gimbal-Lock, zu vermeiden und eine nahtlose Integration in Frameworks wie ROS 2 und MoveIt sicherzustellen.
        Da das traditionelle G-Code-Format primär auf drei Achsen ausgelegt ist und Orientierungsinformationen nur unzureichend abbilden kann, ist die Erzeugung eines formatunabhängigen Datenmodells, das im Projekt als \gls{json}-Datenartefakt umgesetzt ist, erforderlich.
        Dieses Format führt Geometrie, Toolpath, Orientierung und maschinenspezifische Parameter durchgängig zusammen.

      \subsubsection{Wie unterscheidet er sich von einem klassischen, planaren Slicer?}
        Der wesentliche Unterschied liegt in der geometrischen Schichtbildung und der Variabilität der Werkzeugachse.
        Ein klassischer, planarer Slicer zerlegt das Bauteil ausschließlich in horizontale Schichten entlang einer statischen, globalen Aufbauachse.
        Im Gegensatz dazu nutzt der 6-Achs-Slicer dynamische Ansätze:
        \begin{itemize}
        \item Schichtgeometrie: Während planare Slicer nur ebene Schichten erzeugen, ermöglicht ein 6-Achs-Slicer multidirektionale Strategien.
        Der Base-Algorithmus teilt das Bauteil in Stamm und Zweige auf, die in unterschiedlichen lokalen Wachstumsrichtungen gesliced werden.
        Der Kristall-Algorithmus hingegen nutzt vollständig gekrümmte Iso-Flächen basierend auf einem Skalarfeld, sodass sich die Druckrichtung kontinuierlich anpasst
        \item Werkzeugorientierung: Bei einem planaren Slicer fällt die Werkzeugachse konstant mit der globalen Aufbauachse zusammen.
        Bei nicht-planaren Strategien variiert die Orientierung stetig; sie wird beispielsweise aus den Oberflächennormalen (beim Base-Algorithmus) oder dem Gradientenfeld des Skalarfeldes berechnet.
        \item Umgang mit Überhängen: In der planaren Fertigung erfordern Überhänge (typischerweise ab einem Winkel von ca.
        45°) den Einsatz von Stützstrukturen, was den Materialverbrauch und die Nachbearbeitung erhöht.
        Ein 6-Achs-Slicer kann durch die freie Orientierung der Düse überhängende Bereiche so beschichten, dass sie sich durch die bereits gedruckte Basis selbst stützen, wodurch eine stützstrukturfreie Fertigung ermöglicht wird.
        \end{itemize}
        Diese Unterschiede sind in \ref{sec:concept-of-slicing-strategies} detailliert gegenübergestellt.
        \sig{Ben Stahl}


    \subsection{Forschungsfrage 2}
      \subsubsection{Wie kann ein solcher Slicer prototypisch umgesetzt werden?}
        Die Umsetzung des Slicers erfolgt im Projekt durch eine klare Aufteilung in zwei funktionale Softwarekomponenten den Modulen \ac{medusa} und \ac{kronos}.
        Dies geschah um die geometrische Planung von der maschinennahen Ausführung zu trennen:
        \begin{itemize}
        \item \textbf{\acs{medusa} (Slicing):} Dieses Modul ist für das Laden des 3D-Modells, die Visualisierung und den eigentlichen Slicing-Prozess zuständig.
        Es berechnet die Toolpaths und exportiert diese als plattformunabhängiges Datenartefakt.
        Hier sind die verschiedenen Slicing-Strategien wie der Base- oder der Kristall-Algorithmus implementiert.
        \item \textbf{\acs{kronos} (Robotersteuerung):} Diese Komponente fungiert als Middleware.
        Sie nimmt den Toolpath aus \ac{medusa} entgegen und transformiert die geometrischen Daten in Gelenktrajektorien für den Roboter, hier ein Universal Robots UR3e.
        Die Ansteuerung erfolgt dabei über die Frameworks ROS 2 und MoveIt 2
        \end{itemize}
        Aufgearbeitet wurde diese Trennung im Abschnitt \ref{sec:overall-system-architecture}.
        
      \subsubsection{Welche technischen Schwierigkeiten treten dabei auf?}
        Bei der Entwicklung und Erprobung traten technische Schwierigkeiten auf, die über das klassische 3D-Drucken hinausgehen:
        \begin{itemize}
        \item \textbf{Komplexität der Bahnplanung und inversen Kinematik:} Im Gegensatz zu kartesischen Druckern muss für jede Position im Raum die Gelenkstellung des Roboters berechnet werden (Inverse Kinematik).
        Dabei können Singularitäten auftreten – Bereiche, die der Roboter mathematisch oder physisch nicht erreichen kann oder in denen Gelenke theoretisch unendliche Geschwindigkeiten erreichen müssten.
        Siehe Abschnitt \ref{subsec:inverse-kinematics-and-singularities}.
        \item \textbf{Fehlende Standardisierung:} Da etablierte Formate wie herkömmlicher G-Code die benötigten 6-\ac{DOF}-Informationen (Position + Orientierung) nicht standardmäßig unterstützen, musste ein eigenes Austauschformat entwickelt werden, um die Datenkette zwischen Slicer und Roboter zu schließen.
        \end{itemize}
        \sig{Ben Stahl}

    \subsection{Forschungsfrage 3}
      \subsubsection{Welche Probleme entstehen bei der Übertragung der geplanten Bahnen auf einen realen Roboter (im konkreten Fall einen UR3e mit ROS und MoveIt)?}
        Die Übertragung der geometrischen Bahndaten auf die Hardware ist mit erheblichen steuerungstechnischen und kinematischen Herausforderungen verbunden, die eine strikte architektonische und algorithmische Behandlung erfordern.
        \begin{itemize}
        \item \textbf{Hohe Wegpunktdichte und Recheneffizienz:} Aufgrund der dichten Wegpunktfolge beim 3D-Druck würde die Einzelverarbeitung von Punkten zu einem zu hohen Kommunikations-Overhead führen.
        Das Middleware-System \ac{kronos} verarbeitet Punkte daher in Clustern von 50 Stück (Batch-Processing), was eine vorausschauende Planung (Lookahead) ermöglicht.
        Siehe \ref{subsubssec:functionality}
        \item \textbf{Bahninterpolation und konstanter Materialfluss:} Für ein sauberes Druckbild darf der Roboter an Wegpunkten nicht exakt stoppen, da jeder Geschwindigkeitsabfall zu einer Materialanhäufung führen würde (vgl. \ref{subsec:path-interpolation-and-path-control}).
        Es muss eine lineare Bahninterpolation (z. B. Cartesian Path Planning mit 0,001 m Schrittweite) angewandt werden, um eine konstante Bahngeschwindigkeit zu erzielen.
        \item \textbf{Kollisionsvermeidung:} Der Roboter bewegt sich in einem engen Arbeitsraum über dem bereits extrudierten Bauteil.
        Eine kontinuierliche Echtzeit Kollisionsprüfung muss erfolgen, um Kollisionen zwischen dem Endeffektor und dem Druckbett oder dem Bauteil zu verhindern.
        \end{itemize}

      \subsubsection{Wie lassen sich die beobachteten Effekte wie Konfigurationswechsel und Annäherungen an Singularitäten einordnen?}
        \textbf{Annäherung an Singularitäten:}
        \begin{itemize}
        \item Singularitäten sind Bereiche, die der Roboterarm aufgrund physischer Gelenklimits nicht erreichen kann oder bei denen der idealisierte mathematische Ansatz zu theoretisch unendlichen Gelenkgeschwindigkeiten führen würde (vgl. \ref{subsec:inverse-kinematics-and-singularities}).
        \item Da der 6-Achs-Druck komplexe Neigungen erfordert, operiert der Roboter oft sehr nah an seinen physikalischen Arbeitsraumgrenzen.
        Normale Solver scheitern hier oft an lokalen Minima.
        Daher wird ein spezialisierter Solver (TRAC-IK) eingesetzt, der iterative und genetische Algorithmen kombiniert, um diese Minima zu überwinden (vgl. \ref{subsec:component-kronos}).
        \item Kann eine Route aufgrund von Singularitäten nicht mit ausreichender Qualität berechnet werden (Coverage unter 95 \%), bricht das System den Prozess ab, um inkonsistente Druckzustände zu vermeiden (vgl. \ref{subsec:component-kronos}).
        \end{itemize}
        \textbf{Konfigurationswechsel (Joint Flipping):}
        \begin{itemize}
        \item Die Inverse Kinematik liefert oft mehrere mögliche Gelenkstellungen für ein und dieselbe Zielposition (vgl. \ref{subsec:inverse-kinematics-and-singularities}).
        \item Ein kritisches Problem entsteht bei kontinuierlichen Bahnen, wenn der Solver zwischen zwei Wegpunkten plötzlich in eine völlig andere Gelenkkonfiguration springt.
        Diese Problematik wird auch „Joint Flipping“ genannt.
        \item Um dies einzuordnen und zu lösen, nutzt das System einen „Unwrapping-Algorithmus“: Überschreitet die Differenz zwischen aufeinanderfolgenden Gelenkwinkeln den Wert von $\pi$, wird durch Addition oder Subtraktion von $2\pi$ sichergestellt, dass das Gelenk nicht den langen, extremen Weg über den Nullpunkt wählt (vgl. \ref{subsec:component-kronos}).
        Ohne diese Korrektur würde der Roboter abrupt umschwenken und die kontinuierliche Extrusion zerreißen.
        \end{itemize}
        \sig{Ben Stahl}

    \subsection{Forschungsfrage 4}
    \subsubsection{Welche Schlussfolgerungen lassen sich aus den umgesetzten und den nicht umgesetzten Slicing-Strategien für eine zukünftige Weiterentwicklung ziehen?}
      Die Arbeit beleuchtet zwei umgesetzte nicht-planare Strategien (den Base- und den Kristall-Algorithmus) sowie eine rein konzeptionell entworfene, nicht umgesetzte Strategie (den Sphären-Algorithmus).
      Aus deren Analyse ergeben sich folgende Handlungsfelder für zukünftige Iterationen:
      \subsubsection{Schlussfolgerungen aus den umgesetzten Strategien (Base- und Kristall-Algorithmus)}

      \paragraph{Integration einer ebenenübergreifenden Kollisionsprüfung (Branch-Kollision)}
        Sowohl beim multidirektionalen Base-Algorithmus als auch beim feldgetriebenen Kristall-Algorithmus wurde festgestellt, dass Zweige (Branches), die aus unterschiedlichen Richtungen aufeinander zu wachsen, derzeit zu überlappenden Schichten führen.
        Eine automatische Verschmelzung findet nicht statt (\ref{subsec:concept-base-algorithm}, \ref{subsec:concept-crystal-algorithm}).
        Eine Weiterentwicklung muss zwingend eine vorausschauende topologische Prüfung integrieren, die Kollisionen zwischen wachsenden Bauteilzweigen erkennt und die Pfade entsprechend verschmilzt oder umleitet.

      \paragraph{Erweiterung der Wandkonturen} 
        Die Erzeugung aufwendiger Füll- und Wandstrukturen auf gekrümmten Flächen ist mathematisch anspruchsvoll.
        Der Kristall-Algorithmus unterstützt aktuell nur eine einzige Wandkontur pro Schicht.
        Um mehrwandige Konturen zu erzeugen, die für die mechanische Stabilität unerlässlich sind, muss in Zukunft ein geodätisches Versatzverfahren (Geodesic Offset) für gekrümmte Iso-Flächen entwickelt werden (\ref{subsec:concept-crystal-algorithm}).
        %ToDo (Optional) Auf Infill-Generiereung noch eingehen 

      \paragraph{Umgang mit komplexen Topologien und Basisregionen} 
        \begin{itemize}
        \item Der Base-Algorithmus beschneidet nach Erkennung eines Überhangs den erlaubten Bereich starr.
        Ein Wechselspiel aus Überhang, vertikalem Stück und erneutem Überhang in derselben Richtung wird dadurch nicht unterstützt (\ref{subsec:concept-base-algorithm}).
        \item Der Kristall-Algorithmus geht zudem von exakt einer zusammenhängenden Buildplate-Kontaktregion aus (\ref{subsec:concept-crystal-algorithm}).
        \item \textbf{Schlussfolgerung:} Zukünftige Algorithmen müssen dynamischer auf mehrteilige Bauplätze oder verschachtelte Überhänge reagieren können, beispielsweise durch das Berechnen separater Skalarfelder pro Teilkomponente.
        \end{itemize}

      \paragraph{Performance-Optimierung der geometrischen Berechnungen} 
        Die feldgetriebenen Ansätze (Kristall-Slicing) sind stark von der Lösung großer linearer Gleichungssysteme abhängig.
        Dies macht das Verfahren so rechenintensiv, dass eine kontinuierliche Echtzeitvorschau in der Nutzeroberfläche derzeit ausgeschlossen ist (\ref{subsec:component-medusa}, \ref{subsec:concept-crystal-algorithm}).
        Für einen produktiven Einsatz muss die Berechnungsgeschwindigkeit massiv erhöht werden.

    \subsubsection{Schlussfolgerungen aus der nicht umgesetzten Strategie (Sphären-Algorithmus)}
      Obwohl der Sphären-Algorithmus nur als Konzept vorliegt, liefert seine theoretische Analyse wertvolle Erkenntnisse für die Konzeption künftiger Strategien:
      
      \paragraph{Kompensation schwacher Schichtübergänge}
        Während die Stabilität innerhalb einer doppelt gekrümmten Sphärenschale theoretisch sehr hoch ist, bilden die Übergänge zwischen aufeinanderfolgenden Sphärenschalen mechanische Schwachstellen, da die Kontaktfläche orthogonal und somit relativ klein ist (\ref{subsec:concept-sphere-algorithm}).
        Eine zukünftige Umsetzung müsste dieses Problem durch angepasste Extrusionsparameter oder ineinandergreifende Schichtübergänge lösen.

      \paragraph{Entwicklung angepasster Infill-Muster für 3D-Flächen}
        Für rein planare Schichten sind Füllmuster leicht generierbar.
        Da Sphärenschalen jedoch doppelt gekrümmt sind, lassen sich 2D-Muster nicht ohne starke Verzerrungen übertragen.
        Eine zukünftige Weiterentwicklung muss dedizierte, auf die Sphäre parametrierte Infill-Verfahren (z. B. spiralförmige Bahnen) definieren (\ref{subsec:concept-sphere-algorithm}).

      \paragraph{Spezialisierung statt Generalisierung (Hybride Ansätze)}
        Der Sphären Algorithmus zeigt, dass Bauteile, deren Hauptmasse nicht radial vom Mittelpunkt erreichbar ist, zu einer ineffizienten Vielzahl an gespawnten Zweigen führen (\ref{subsec:concept-sphere-algorithm}).
        Umgekehrt entartet der Kristall-Algorithmus bei simplen Quadern zu nutzlos aufwendigen planaren Schichten (\ref{subsec:concept-crystal-algorithm}).

      Schlussfolgerung: Zukünftige Slicer-Software sollte sich nicht auf einen einzigen „perfekten“ Algorithmus verlassen.
      Vielmehr bedarf es hybrider Ansätze, bei denen das Bauteil vorab segmentiert wird und für jede Teilgeometrie (radial, organisch-gekrümmt, streng prismatisch) die mathematisch und fertigungstechnisch effizienteste Slicing-Strategie automatisch ausgewählt wird.
      \sig{Ben Stahl}

  \section{Implikationen für Forschung und Praxis}\label{sec:implications-for-research-and-practice}
    Die konsequente Fortführung dieses Projektes birgt ein erhebliches Potenzial, sowohl für die akademische Forschung, als auch für zukunftsweisende industrielle Anwendungen im Bereich der additiven Fertigung.

    \subsection{Industrielle Praxis und wirtschaftliche Relevanz}\label{subsec:industrial-practice and-economic-relevance}
      In der industriellen Praxis ermöglicht die Weiterentwicklung nicht-planarer Slicing-Verfahren den technologischen Übergang zu einem weitestgehend stützstrukturfreien, oder zumindest signifikant stützreduzierten 3D-Druck.
      Durch die Realisierung solcher adaptiven Fertigungsstrategien lassen sich komplexe Bauteile mit einer erheblichen Materialersparnis produzieren, was unmittelbar zur Steigerung der Ressourceneffizienz und zur Senkung der operativen Produktionskosten beiträgt.
      Über die reine Materialeinsparung hinaus führt die gezielte Reduktion von Stützkonstruktionen zu einer drastischen Verringerung des manuellen Nachbearbeitungsaufwands.
      Zeitintensive Prozesse, wie das mechanische Entfernen von Support-Strukturen sowie die anschließende Oberflächenveredelung an den Kontaktstellen, entfallen weitestgehend.
      Dies führt dazu, dass die gesamte Prozesskette ökonomischer, schneller und qualitativ hochwertiger gestaltet werden kann.
      \sig{Ben Stahl}

    \subsection{Wissenschaftliche Forschung und methodische Flexibilität}\label{subsec:scientific-research-and-methodical-flexibility}
      Für die wissenschaftliche Forschung bietet der gewählte modulare Systemaufbau eine ideale und hochflexible Plattform für künftige Untersuchungen.
      Diese Architektur erlaubt es, neuartige Slicing-Algorithmen in einer isolierten Umgebung zu implementieren, experimentell zu testen und systematisch zu validieren, ohne dass eine zeitintensive Neukonfiguration des Gesamtsystems erforderlich ist.
      Durch diese Effizienzsteigerung innerhalb der Entwicklungsumgebung können innovative Ansätze zur multidimensionalen Trajektorienplanung schneller iteriert und in Richtung einer industriellen Anwendungsreife geführt werden.
      Somit fungiert das vorliegende Projekt als katalytische Basis für weiterführende Studien im Bereich der robotergestützten, adaptiven Fertigungsverfahren und setzt neue Impulse für die computergestützte Produktionsplanung.
      \sig{Ben Stahl}